\section*{Chapitre \ref{ch:b2:enolates-aldol} --- Énols, énolates et aldolisation}

\begin{solution}{exo:b2:enolates-aldol:1}
Butanone~: but-1-én-2-ol \ce{CH2=C(OH)CH2CH3} et but-2-én-2-ol \ce{CH3C(OH)=CHCH3} ((\textit{E}) et
(\textit{Z})). Cyclohexanone~: cyclohex-1-én-1-ol.
\end{solution}

\begin{solution}{exo:b2:enolates-aldol:2}
Éthane $<$ propanone $<$ 3-oxobutanoate d'éthyle (10.7) $<$ pentane-2,4-dione (9.0).
\end{solution}

\begin{solution}{exo:b2:enolates-aldol:3}
Le 3-hydroxybutanal, puis le but-2-énal par chauffage.
\end{solution}

\begin{solution}{exo:b2:enolates-aldol:4}
Propanedioate de diéthyle, éthanolate de sodium, 1-bromobutane~; hydrolyse en acide butylpropanedioïque~;
le chauffage élimine \ce{CO2}~: \ce{CH3(CH2)3CH2COOH}, l'acide hexanoïque.
\end{solution}

\begin{solution}{exo:b2:enolates-aldol:5}
Cinétique~: déprotonation sur C6 (du côté du \ce{CH2}), par le LDA à \qty{-78}{\degreeCelsius}. Thermodynamique~:
la double liaison la plus substituée C1=C2, vers le carbone qui porte le méthyle, avec un alcoolate dans
son alcool à température ambiante.
\end{solution}

\begin{solution}{exo:b2:enolates-aldol:6}
On obtient la 4-phénylbut-3-én-2-one (benzalacétone), \ce{C6H5CH=CHCOCH3}. Le benzaldéhyde n'a pas d'hydrogène $\alpha$~:
il ne peut pas former d'\omterm{def:b2:enolates-aldol:enolate}{énolate} et n'est que l'électrophile.
\end{solution}

\begin{solution}{exo:b2:enolates-aldol:7}
L'\omterm{def:b2:enolates-aldol:enolate}{énolate} de l'un des groupes méthyle (C1) attaque l'autre carbonyle (C5)~: un cycle à cinq atomes,
la 3-méthylcyclopent-2-én-1-one après déshydratation. L'autre possibilité, l'\omterm{def:b2:enolates-aldol:enolate}{énolate} en C3 attaquant C5,
formerait un cycle tendu à trois atomes.
\end{solution}

\begin{solution}{exo:b2:enolates-aldol:8}
\ce{2CH3COOC2H5 -> CH3COCH2COOC2H5 + C2H5OH} par l'éthanolate de sodium~; le produit, déprotoné, est
récupéré sous forme de 3-oxobutanoate d'éthyle par acidification.
\end{solution}

\begin{solution}{exo:b2:enolates-aldol:9}
La propanone, l'éthanal et l'acétophénone (ils possèdent un groupe \ce{CH3-CO}~; l'éthanal est \ce{CH3CHO}).
La pentan-3-one, non.
\end{solution}

\begin{solution}{exo:b2:enolates-aldol:10}
L'\omterm{def:b2:enolates-aldol:enolate}{énolate} formé à une extrémité attaque l'ester de l'autre~: le 2-oxocyclopentane-1-carboxylate d'éthyle,
un cycle à cinq atomes.
\end{solution}

\begin{solution}{exo:b2:enolates-aldol:11}
Dans la 3-méthylpentan-2-one, le stéréocentre (C3) est le carbone $\alpha$~: l'énolisation le rend plan,
et la reprotonation par l'une ou l'autre face le racémise. Dans la 4-méthylhexan-2-one, le stéréocentre
est en C4, en $\beta$ du carbonyle, et l'énolisation ne le touche pas.
\end{solution}

\begin{solution}{exo:b2:enolates-aldol:12}
La propanone avec deux équivalents de benzaldéhyde en milieu basique. Avec deux aldéhydes différents,
condenser d'abord la propanone avec l'un des aldéhydes (un équivalent, en isolant la benzalacétone), puis
condenser son groupe méthyle restant avec le second. En une seule étape, les deux aldéhydes entrent en
compétition et donnent un mélange de trois diénones.
\end{solution}

\begin{solution}{pb:b2:enolates-aldol:1}
\textbf{1.} La propanone, qui porte six hydrogènes $\alpha$, trois sur chacun de ses groupes méthyle.
\textbf{2.} $\mathrm{CH_3{-}CO{-}CH_2^-} \leftrightarrow \mathrm{CH_3{-}C(O^-){=}CH_2}$.
\textbf{3.} Il n'a pas d'hydrogène sur le carbone voisin de son carbonyle (ce carbone appartient au cycle
et ne porte pas de H).
\textbf{4.} Le benzaldéhyde, un aldéhyde en excès, est le meilleur électrophile~; l'\omterm{def:b2:enolates-aldol:enolate}{énolate} de la propanone
le rencontre plus souvent qu'un carbonyle de propanone.
\textbf{5.} Une cétone porte deux groupes carbonés qui repoussent la densité électronique vers le carbone
de son carbonyle et gênent l'approche~; un aldéhyde n'en porte qu'un.
\textbf{6.} Seule une petite fraction d'\omterm{def:b2:enolates-aldol:enolate}{énolate} est nécessaire~: il est régénéré à mesure qu'il réagit.
\textbf{7.} Le carbone de l'\omterm{def:b2:enolates-aldol:enolate}{énolate} se lie au carbone du carbonyle du benzaldéhyde~: on obtient l'alcoolate de
la 4-hydroxy-4-phénylbutan-2-one.
\textbf{8.} Arrachement d'un hydrogène $\alpha$ et départ de l'hydroxyde du carbone $\beta$~: la 4-phénylbut-3-én-2-one.
\textbf{9.} La double liaison formée est conjuguée à la fois avec le cycle et avec le carbonyle.
\textbf{10.} L'autre groupe méthyle possède encore trois hydrogènes $\alpha$.
\textbf{11.} \ce{2C6H5CHO + CH3COCH3 -> C6H5CH=CHCOCH=CHC6H5 + 2H2O}.
\textbf{12.} Les déshydratations, et la précipitation du produit, qui quitte la solution.
\textbf{13.} Deux, toutes deux (\textit{E}).
\textbf{14.} Les isomères (\textit{Z}) placent le groupe phényle près du groupe carbonyle~: gêne stérique.
\textbf{15.} Deux \ce{C=C}, un \ce{C=O} et les deux cycles~: un système conjugué qui traverse la molécule.
\textbf{16.} Sa conjugaison étendue réduit l'écart $\pi$ (\cref{prop:b2:huckel:gap})~: elle absorbe à
l'extrémité bleue du visible et paraît jaune.
\textbf{17.} Non~: elle est plane (ou presque) et possède un plan de symétrie.
\textbf{18.} 106.12, 58.08 et \qty{234.30}{g/mol}.
\textbf{19.} Benzaldéhyde~: $2.1 \times 1.045/106.12 = \qty{20.7}{mmol}$~; propanone~: $0.75 \times
0.7845/58.08 = \qty{10.1}{mmol}$.
\textbf{20.} La propanone exige \qty{20.3}{mmol} de benzaldéhyde, et l'on en dispose de \qty{20.7}{mmol}~:
la propanone est limitante, de justesse.
\textbf{21.} $0.01013 \times 234.30 = \qty{2.37}{g}$.
\textbf{22.} Le benzaldéhyde restant demeure dans l'éthanol et est éliminé au lavage.
\textbf{23.} Le produit de monocondensation, la benzalacétone.
\textbf{24.} $0.75 \times 2.37 = \boldsymbol{\approx \qty{1.78}{g}}$ de dibenzalacétone.
\end{solution}
